\clearpage\begingroup\SourceLeaf{67b}\HistoricalHead{62}{L}{PHYSIK. NACHLASS.}
\begin{center}
\rule{.25\textwidth}{.35pt}\\
{\small [Aus Handbuch 15, Ba.]}\\
\rule{.25\textwidth}{.35pt}\\[.8em]
[4.]
\end{center}
\noindent{\small [S. 191]}
\begin{center}\textbf{Stärke eines Inductionsstosses.}\end{center}
\noindent\hspace*{\parindent}Es sei \(a\) die Elongation, auf welche die Nadel getrieben werden würde,
wenn keine Dämpfung vorhanden wäre. \(\varphi,\theta\) in der Bedeutung der Res. 1837[\textsuperscript{*})]\begingroup\renewcommand{\thefootnote}{}\footnotetext{[\textsuperscript{*}) \textit{Anleitung zur Bestimmung der Schwingungsdauer einer Magnetnadel}, Resultate u.s.w. 1837, S. 58; Werke V, S. 374.]}\endgroup.

Es wird dann
\begin{center}
\renewcommand{\arraystretch}{2.5}
\begin{tabular}{ll}
erste wirkliche Elongation&\( +a\theta^{\frac12-\frac{\varphi}{\pi}}\)[\textsuperscript{**})],\\
zweite&\( -a\theta^{\frac32-\frac{\varphi}{\pi}}\),\\
dritte&\( +a\theta^{\frac52-\frac{\varphi}{\pi}}\)
\end{tabular}\\
u.s.w.\par\medskip
\begin{tabular}{ll}
also erster Schwingungsbogen&\(a(1+\theta)\theta^{\frac12-\frac{\varphi}{\pi}}\),\\
zweiter&\(a(1+\theta)\theta^{\frac32-\frac{\varphi}{\pi}}\).
\end{tabular}
\end{center}
\begingroup\renewcommand{\thefootnote}{}\footnotetext{[\textsuperscript{**}) In der Handschrift steht hier und in den folgenden Formeln an entsprechender Stelle noch der Faktor \(\cos\varphi\).]}\endgroup
\clearpage\endgroup
